\section{Comprensión de los datos}

\subsection{Fuente}
\url{https://www.kaggle.com/datasets/erdemtaha/cancer-data}

\subsection{Descripción}
El dataset contiene mediciones morfológicas de núcleos celulares obtenidas de imágenes de masas mamarias (biopsia por aspiración con aguja fina). El objetivo es distinguir entre tumores benignos y malignos.

\subsection{Variables del dataset}
\begin{longtable}{p{0.3\textwidth} p{0.7\textwidth}}
\toprule    
\textbf{Variable} & \textbf{Descripción} \\
\midrule
\endhead
\texttt{id} & Identificador único del caso/paciente. \\
\texttt{diagnosis} & Diagnóstico (B = Benigno, M = Maligno). \\
\texttt{radius\_mean} & Promedio del radio. \\
\texttt{texture\_mean} & Promedio de textura (variación de niveles de gris). \\
\texttt{perimeter\_mean} & Promedio del perímetro. \\
\texttt{area\_mean} & Promedio del área. \\
\texttt{smoothness\_mean} & Promedio de suavidad. \\
\texttt{compactness\_mean} & Promedio de compacidad. \\
\texttt{concavity\_mean} & Promedio de concavidad. \\
\texttt{concave points\_mean} & Promedio de puntos cóncavos. \\
\texttt{symmetry\_mean} & Promedio de simetría. \\
\texttt{fractal\_dimension\_mean} & Promedio de dimensión fractal. \\
\texttt{radius\_se} & Error estándar del radio. \\
\texttt{texture\_se} & Error estándar de textura. \\
\texttt{perimeter\_se} & Error estándar del perímetro. \\
\texttt{area\_se} & Error estándar del área. \\
\texttt{smoothness\_se} & Error estándar de suavidad. \\
\texttt{compactness\_se} & Error estándar de compacidad. \\
\texttt{concavity\_se} & Error estándar de concavidad. \\
\texttt{concave points\_se} & Error estándar de puntos cóncavos. \\
\texttt{symmetry\_se} & Error estándar de simetría. \\
\texttt{fractal\_dimension\_se} & Error estándar de dimensión fractal. \\
\texttt{radius\_worst} & Valor peor/máximo del radio. \\
\texttt{texture\_worst} & Valor peor/máximo de textura. \\
\texttt{perimeter\_worst} & Valor peor/máximo del perímetro. \\
\texttt{area\_worst} & Valor peor/máximo del área. \\
\texttt{smoothness\_worst} & Valor peor/máximo de suavidad. \\
\texttt{compactness\_worst} & Valor peor/máximo de compacidad. \\
\texttt{concavity\_worst} & Valor peor/máximo de concavidad. \\
\texttt{concave points\_worst} & Valor peor/máximo de puntos cóncavos. \\
\texttt{symmetry\_worst} & Valor peor/máximo de simetría. \\
\texttt{fractal\_dimension\_worst} & Valor peor/máximo de dimensión fractal. \\
\bottomrule
\end{longtable}

\subsection{Tipos de variables, faltantes y tamaño}
\begin{itemize}
    \item Variables categóricas: 1 (\texttt{diagnosis}).
    \item Variables numéricas: 31 (\texttt{id} + 30 mediciones clínicas continuas).
    \item Faltantes detectados inicialmente: una columna extra sin nombre \texttt{Unnamed: 32}, completamente vacía.
    \item Tamaño del dataset tras limpieza: 569 filas y 32 columnas.
\end{itemize}

\subsection{Preguntas guía}
\begin{itemize}
    \item \textbf{Variables más relevantes:} medidas de tamaño, forma e irregularidad del tumor (por ejemplo \texttt{radius\_mean}, \texttt{perimeter\_mean}, \texttt{area\_mean}, \texttt{radius\_worst}, \texttt{concave points\_worst}).
    \item \textbf{Variables redundantes o inútiles:} pueden existir variables que aportan información similar entre sus versiones \texttt{mean}, \texttt{se} y \texttt{worst}; además \texttt{id} no se usa para modelar el fenómeno clínico.
    \item \textbf{Sesgos:} hay más casos benignos que malignos, lo cual implica desbalance de clases.
    \item \textbf{Calidad de datos:} se encontró una columna vacía por formato del archivo y se removió en preparación.
\end{itemize}